% ---------------------------------------------------------------------------
%  Section 5 -- Conclusion
%
%  TODO once the full test sweep lands: insert the headline numbers in the
%  first paragraph, marked XX below. Everything else is result-independent.
% ---------------------------------------------------------------------------

\section{Conclusion}
\label{sec:conclusion}

We presented a three-dimensional conditional latent diffusion framework for
whole-tumour segmentation from incomplete multi-sequence MRI. Two volumetric
autoencoders compress images and annotations by a factor of four per axis,
which makes diffusion over $128^{3}$ volumes tractable, and modality dropout
during training lets a single model serve all fifteen acquisition
configurations. Repeated sampling yields both a consensus segmentation and a
voxel-wise map of disagreement between plausible annotations.

Beyond accuracy, we treated modality availability as an experimental variable
rather than a nuisance. Because withholding a sequence cannot add information
about the annotation, a calibrated model should report more uncertainty as
sequences are removed, and evaluating every case under every subset makes that
prediction testable within each case. Our results indicate that the predicted
uncertainty follows this behaviour, which suggests it reflects genuine
information loss rather than sampling noise. % TODO: quote the correlation

% ---------------------------------------------------------------------------
\subsection{Limitations and Future Work}
\label{subsec:limitations}

\paragraph{Compression cost on small structures.}
The autoencoders apply a fixed $4\times$ downsampling per axis, chosen so that
diffusion over whole volumes fits within memory. This is not neutral with
respect to structure size. Reconstructing the annotation through the frozen
autoencoder recovers the whole tumour almost exactly, but is markedly less
faithful for the enhancing core, which is smaller and thinner. Since the
diffusion model operates entirely inside this latent space, that gap is a
ceiling it cannot exceed, and it is the main reason we report the whole tumour
rather than the full BraTS region hierarchy. Adaptive or anisotropic
downsampling, or a latent with higher spatial resolution for the annotation
branch, would target this directly.

\paragraph{Region coverage.}
We segment the whole tumour only. Extending to tumour core and enhancing
tumour is the natural next step, and the analysis above suggests it depends
more on the compression scheme than on the diffusion model itself. Confirming
that would require repeating the autoencoder study at several downsampling
factors and measuring where the reconstruction ceiling for each region falls.

\paragraph{Reconstruction objective.}
The annotation autoencoder combines weighted cross-entropy with a soft Dice
term. Losses designed specifically for thin or sparse structures, such as
boundary-aware or focal variants, may preserve small regions better at the
same compression factor, and we have not compared them.

\paragraph{Uncertainty baselines.}
We compare the predicted uncertainty against the behaviour information theory
prescribes, but not against other estimators. Monte Carlo dropout and deep
ensembles would place our estimates on a common footing with established
methods, and the Probabilistic U-Net would provide a latent-variable
comparison. Such a study would also allow calibration metrics that require a
reference distribution, which a single fused annotation per case does not
support.
